% page 450 of the facsimile (image p-449 of the scan)
% block 152.4 x 233.3 mm ; left margin 8.0 mm
\pgc{-12.700mm}{0.000mm}{
\l{\cel{3}442}
\l{}
\l{\sou{pistar}. (\sou{trans}.)Reduktar a mikra fragmenti, per shoki}
\l{\cel{3}iterata. - EFI.}
\l{}
\l{\sou{pistacho}. (\sou{bot.)}\cel{2}Mandelo de frukto di arboro ek la fa-}
\l{\cel{3}milio \textquotedbl{}terebintacei\textquotedbl{}, envolvita en pelikulo verda, uza-}
\l{\cel{3}ta da la sukrajisti. - L.\cel{1}\sou{pistacia vera}. - DEFIRS.}
\l{}
\l{\sou{pistolo}. Pafarmo mikra, tre portebla, quan onu tenas en}
\l{\cel{3}nur un ek la manui kande onu pafas. - DEFIRS.}
\l{}
\l{\sou{pistono}. (\sou{tekn.}) Cilindro,\cel{2}movanta en pumpilo per stango,}
\l{\cel{3}retropulsante fluido o pulsata da lu. - DEFIRS.}
\l{}
\l{\sou{pitagorala}. I. (\sou{filoz}.) Apartenanta a Pithagores, a lua}
\l{\cel{3}penso-skolo od a lua doktrini. - II. (\sou{matem.}) Nombri pi-}
\l{\cel{3}tagorala : nombri interligita per la relato\cel{2}a2 = b2 + c}
\l{}
\l{\sou{pitekantropo}. (\sou{paleont.}) Genero de granda simii fosila,}
\l{\cel{3}reprezentata da kranio-kaloto, kelka denti ed un femuro,}
\l{\cel{3}trovita en la yaro l89l en sulo-strato pliocena o pleis-}
\l{\cel{3}tocena, en la insulo Java (Azia). - DEFIRS.}
\l{}
\l{\sou{pitio}. (\sou{epoki antiqua}). Sacerdotino di la oraklo di Apolono,}
\l{\cel{3}en Delphos. - DEFIS.}
\l{}
\l{\sou{pitono}.\cel{1}\sou{(zool.}) Sorto di boao ek India. - L.\cel{1}\sou{python}.-}
\l{\cel{3}II. (\sou{mitol}.) Nomo di serpento monstra, mortigita da}
\l{\cel{3}Apolono, apud Delphos. - DEFIS.}
\l{}
\l{\sou{pitoniso}. (\sou{biblo}).Muliero qua anuncas la eventontaji. -}
\l{\cel{3}DEFIS.}
\l{}
\l{\sou{pitoreska}. I. Qua, pro lua dispozeso, aspektas apta dive-}
\l{\cel{3}nar piktajo. - II. Dicesas pri stilo qua quaze piktas}
\l{\cel{3}la personi o la kozi. - DEFIS.}
\l{}
\l{\sou{pituito}. (\sou{patol.}) Humoro viskoza sekrecata da ula mukosi}
\l{\cel{3}di la nazo, di la bronkii, en stando morbala. -EFIS.}
\l{}
\l{\sou{pivoto}. - (\sou{mekan.}) Suportilo di la axo cirkum qua rotacas}
\l{\cel{3}korpo. - II. (milit-arto). (aludante konversiono di}
\l{\cel{3}\sou{trupo de soldati}).Punto cirkum qua ta konversiono even-}
\l{\cel{3}tas. - III. (\sou{metaf}.) Motivo precipua di la agi, che la}
\l{\cel{3}homi. - EF.}
\l{}
\l{\sou{pizo}.\cel{2}(\sou{bot.}) Singla de la grani de planto leguminosa,}
\l{\cel{3}papilionacea (L.\cel{1}\sou{pisus sativus}), inkluzita en shelo}
\l{\cel{3}verda. - EFI.}
\l{}
\l{\sou{placo}. Loko publika, netektoza, perimetro-garnisata per do-}
\l{\cel{3}mi o monumenti, kun, posible, monumento centrala. -}
\l{\cel{3}- DEFIRS.}
}
